% from ega4-1-fr.tex
\oldpage[IV-1]{233}

de présentation finie~; donc (1.4.5), si $\mathfrak J$ est le noyau de $u$,
le $\widetilde{B'}$-Module $\widetilde{\mathfrak J}$ est de type fini, et par
suite (II, 6.1.4.1) $\mathfrak J$ est un idéal de type fini.
\end{proof}

\begin{proposition}[1.4.7]
\label{IV.1.4.7-fr}
Soit $\rho:A\to B$ un homomorphisme d'anneaux faisant de $B$ une $A$-algèbre
finie. Pour que $B$ soit une $A$-algèbre de présentation finie, il faut et il
suffit que $B$ soit un $A$-module de présentation finie.
\end{proposition}

Nous utiliserons le

\begin{lemma}[1.4.7.1]
\label{IV.1.4.7.1-fr}
Si $B$ est une $A$-algèbre finie, il existe un $A$-homomorphisme surjectif
d'algèbres $u:B'\to B$, où $B'$ est une $A$-algèbre finie et de présentation
finie qui est un $A$-module libre.
\end{lemma}

\begin{proof}
En effet, il y a un système fini $(a_i)_{1\leq i\leq m}$ de générateurs de la
$A$-algèbre $B$ tel que chaque $a_i$ vérifie une relation $F_i(a_i)=0$, où
$F_i\in A[X]$ est un polynôme unitaire de degré $>0$~; la $A$-algèbre quotient
$B'_i=A[X]/(F_i)$ est donc libre de rang fini sur $A$~; soit $c_i$ la classe de
$X$ dans $B'_i$. Il existe un $A$-homomorphisme d'algèbres $u_i:B'_i\to B$
tel que $u_i(c_i)=a_i$~; il suffit alors de prendre pour $B'$ le produit
tensoriel $B'_1\otimes_AB'_2\otimes_A\cdots\otimes_AB'_m$ et de considérer
l'homomorphisme $u_1\otimes u_2\otimes\cdots\otimes u_m:B'\to B$, qui est
surjectif par construction. En outre, si $B''=A[T_1,\ldots,T_m]$ et si $b'_i$
désigne l'image canonique de $c_i$ dans $B'$, le $A$-homomorphisme d'algèbres
$v:B''\to B'$ tel que $v(T_i)=b'_i$ ($1\leq i\leq m$) est surjectif, et son
noyau $\mathfrak b'$ est l'idéal de $B''$ engendré par les $F_i(T_i)$, donc de
type fini~; ceci prouve que $B'$ est une $A$-algèbre de présentation finie.
\end{proof}

\begin{proof}
Ce lemme étant démontré, gardons les mêmes notations, et soit
$w=v\circ u:B''\to B$. Si $\mathfrak b$ (resp. $\mathfrak a$) est le noyau de
$w$ (resp. $u$), on a $\mathfrak a=v(\mathfrak b)$ puisque $v$ est surjectif.
Or (1.4.4), si $B$ est une $A$-algèbre de présentation finie, $\mathfrak b$
est un idéal de type fini de $B''$, donc $\mathfrak a$ est un idéal de type
fini de $B'$, donc un $A$-module de type fini puisque $B'$ est une $A$-algèbre
finie~; comme $B'$ est un $A$-module libre, $B$ est un $A$-module de
présentation finie. Inversement, si $B'$ est un $A$-module de présentation
finie, $\mathfrak a$ est un $A$-module de type fini (Bourbaki,
\emph{Alg. comm.}, chap. I, \S~2, n\textsuperscript{o}~8, lemme 9) et
\emph{a fortiori} un idéal de type fini de $B'$~; par suite, $B$ est par
définition une $B'$-algèbre de présentation finie, et comme $B'$ est une
$A$-algèbre de présentation finie, $B$ est une $A$-algèbre de présentation
finie.
\end{proof}

\subsection{Morphismes de type fini}
\label{subsection:IV.1.5-fr}

\begin{proposition}[1.5.1]
\label{IV.1.5.1-fr}
Soient $f:X\to Y$ un morphisme, $(U_\alpha)$ un recouvrement de $Y$ formé
d'ouverts affines. Les conditions suivantes sont équivalentes~:
\begin{enumerate}
  \item[a)] $f$ est localement de type fini et quasi-compact.
  \item[b)] Pour tout $\alpha$, $f^{-1}(U_\alpha)$ est réunion finie d'ouverts
  affines $V_{\alpha i}$ tels que chaque anneau $A(V_{\alpha i})$ soit une
  $A(U_\alpha)$-algèbre de type fini.
  \item[c)] Pour tout ouvert affine $U$ de $Y$, $f^{-1}(U)$ est réunion finie
  d'ouverts affines $V_j$ tels que les anneaux $A(V_j)$ soient des
  $A(U)$-algèbres de type fini.
\end{enumerate}
\end{proposition}

Pour la démonstration, voir (I, 6.3.2 et 6.6.3).

\begin{definition}[1.5.2]
\label{IV.1.5.2-fr}
On dit qu'un morphisme $f$ est \emph{de type fini} s'il vérifie les conditions
équivalentes de la proposition (1.5.1).
\end{definition}

